\section{Collective coupling capacity and its exact regularization}
\label{sec:collective-capacity}

This section concerns fixed actual marginals and a fixed pure target, not
optimization over local target orientations. The distinction matters:
even commuting, diagonal marginals produce a strict advantage when several
copies are optimized collectively. The calculation reduces the quantum
optimization to the classical Chernoff overlap of two generally
unnormalized measures. The Chernoff mechanism itself is established prior
art; the result here identifies its exact role in this coupling problem.

Let $a,b,p$ be strictly positive probability vectors on $\{1,\ldots,d\}$,
and put $A=\diag a$, $B=\diag b$ and
$x_p=\sum_i\sqrt{p_i}|ii\rangle$. Define
\[
 C_n(a,b;p)=\max_{\rho_A=A^{\otimes n},\ \rho_B=B^{\otimes n}}
       \langle x_p^{\otimes n}|\rho|x_p^{\otimes n}\rangle.
\]
The tensor factors are grouped across the bipartition
$A_1\cdots A_n:B_1\cdots B_n$. Set
\[
 u_i=\sqrt{p_i a_i},\qquad v_i=\sqrt{p_i b_i},\qquad
 Z(t)=\sum_i u_i^t v_i^{1-t},\quad 0\le t\le1,
 \qquad \kappa=\min_{0\le t\le1}Z(t).
\]

\begin{theorem}[Exact collective formula, strict gain, and asymptotic rate]
\label{thm:collective-capacity}
For every positive integer $n$,
\begin{equation}\label{eq:collective-exact}
 C_n(a,b;p)=\left[\sum_{\boldsymbol i\in\{1,\ldots,d\}^n}
          \min\left(\prod_{j=1}^n u_{i_j},
                     \prod_{j=1}^n v_{i_j}\right)\right]^2.
\end{equation}
The maximum is attained by the coherent-minimum coupling of the product
populations $a^{\otimes n},b^{\otimes n}$. If $a\ne b$, then
\begin{equation}\label{eq:collective-strict}
 C_{n+m}>C_n C_m \quad(n,m\ge1),\qquad
 \kappa^2>C_1,
\end{equation}
and
\begin{equation}\label{eq:collective-rate}
 \boxed{\displaystyle
 \lim_{n\to\infty} C_n^{1/n}
   =\left[\min_{0\le t\le1}
       \sum_i\sqrt{p_i}\,a_i^{t/2}b_i^{(1-t)/2}\right]^2
   =\kappa^2.}
\end{equation}
If $a=b$, instead $C_n=C_1^n$ for every $n$.
\end{theorem}

\begin{proof}
The exact diagonal-marginal theorem applied to the product populations
gives
\[
 C_n=\left[\sum_{\boldsymbol i}
       \sqrt{p_{\boldsymbol i}\min(a_{\boldsymbol i},
                                     b_{\boldsymbol i})}\right]^2,
\]
where $a_{\boldsymbol i}=\prod_j a_{i_j}$, and similarly for $b,p$.
Each summand is exactly the minimum in
\eqref{eq:collective-exact}. This also proves attainment by an actual
positive trace-one matrix, not merely a feasible scalar vector.

For nonnegative numbers $U,V,U',V'$,
\[
 \min(UU',VV')\ge\min(U,V)\min(U',V').
\]
If $U>V>0$ and $V'>U'>0$, this inequality is strict.
When $a\ne b$, normalization supplies indices $i,j$ with $a_i>b_i$
and $a_j<b_j$. The all-$i$ string of length $n$ and all-$j$ string of
length $m$ give a strict term. Summing the nonnegative differences
and squaring proves \eqref{eq:collective-strict}'s first assertion.
Also $u_i^t v_i^{1-t}\ge\min(u_i,v_i)$ for every $t$.
For $0<t<1$ at least one term is strict, and both endpoints have a
strict term because the two signs of $a_i-b_i$ occur. Continuity on
$[0,1]$ then yields $\kappa>\sum_i\min(u_i,v_i)=\sqrt{C_1}$.

For the upper asymptotic estimate, apply
$\min(U,V)\le U^tV^{1-t}$ to each product string. The sum factorizes,
so $C_n\le\kappa^{2n}$. For completeness, an elementary matching
lower estimate is established next; it proves the limit without invoking
an unproved statistical asymptotic theorem.
\end{proof}

\begin{theorem}[Finite-copy lower certificate]
\label{thm:finite-copy}
Choose a minimizer $t_*$ of $Z$ and define
\[
 q_i=\frac{u_i^{t_*}v_i^{1-t_*}}{\kappa},\qquad
 \ell_i=\log(u_i/v_i),\qquad
 \mu=\sum_iq_i\ell_i,\qquad
 \sigma^2=\sum_iq_i(\ell_i-\mu)^2.
\]
For every $n\ge1$ and $0<\eta<1$,
\begin{equation}\label{eq:chernoff-sandwich}
 (1-\eta)^2\exp\left(-2\sigma\sqrt{n/\eta}\right)\kappa^{2n}
 \ \le\ C_n\ \le\ \kappa^{2n}.
\end{equation}
If $\sigma=0$, equality $C_n=\kappa^{2n}$ holds.
\end{theorem}

\begin{proof}
For $n$ independent $q$-distributed letters let $Y_n=\sum_j\ell_{i_j}$.
Factoring the products in \eqref{eq:collective-exact} gives the exact identity
\[
 \sqrt{C_n}=\kappa^n\,
 \mathbb E_q\exp\bigl(\min((1-t_*)Y_n,-t_*Y_n)\bigr).
\]
If $0<t_*<1$, differentiation of $Z$ gives $\mu=0$.
At $t_*=0$, the right derivative is nonnegative and $\mu\ge0$;
at $t_*=1$, the left derivative is nonpositive and $\mu\le0$.
For $\sigma>0$, Chebyshev's inequality gives probability at least
$1-\eta$ to the event
$|Y_n-n\mu|<\sigma\sqrt{n/\eta}$, with non-strict endpoints usable
without changing the bound. In the interior case $|Y_n|$ is bounded by
this threshold. At $t_*=0$ only the lower bound on $Y_n$ matters, and
$n\mu\ge0$; at $t_*=1$ only the upper bound matters, and $n\mu\le0$.
In all three cases the exponent in the expectation is at least
$-\sigma\sqrt{n/\eta}$ on this event. Squaring proves
\eqref{eq:chernoff-sandwich}. If $\sigma=0$, all $\ell_i$ coincide.
The minimizing endpoint, or either endpoint in the constant-zero case,
then gives the exact equality directly. Taking logarithms of the two
bounds and dividing by $n$ proves \eqref{eq:collective-rate}.
\end{proof}

The probability normalization of $q$ is essential: $u$ and $v$ themselves
usually have total mass below one. Replacing them silently by normalized
probability vectors changes the minimization unless their total masses
coincide. In classical language, $\sum\min(u^{\otimes n},v^{\otimes n})$
is an overlap for unnormalized product measures. The proof above explicitly
retains their masses.

\paragraph{Zero populations.}
The same exact formula holds with zero entries. For the regularization,
restrict to $I=\{i:p_i a_i b_i>0\}$ before defining $u,v,Z$.
Every surviving string lies in $I^n$. If $I$ is empty, all capacities
vanish. If it is nonempty, use the positive restricted weights without
renormalizing them. The strict-supermultiplicativity criterion is precisely
the existence in $I$ of both $a_i>b_i$ and $a_j<b_j$.
When all active ratios have one sign, the same side is the minimum in
every product and $C_n=C_1^n$. Thus the criterion is necessary as well
as sufficient for the nonnegative case.

\subsection{An exact two-qubit counterexample to multiplicativity}
Take
\[
 A=\diag(7/10,3/10),\qquad B=\diag(3/10,7/10),
 \qquad x=\frac{|00\rangle+|11\rangle}{\sqrt2}.
\]
The one-copy optimum is $C_1=3/5$, attained by the state already used
to rule out optimization over normalized pure couplings. Two copies have
the diagonal minima
\[
 (9/100,21/100,21/100,9/100),
\]
and therefore
\begin{equation}\label{eq:capacity-counterexample}
 \boxed{C_2=\frac{15+3\sqrt{21}}{50}
       \simeq0.57495454\ >\ \frac9{25}=C_1^2.}
\end{equation}
An exact attaining four-by-four bipartite state is
\[
 \rho_2=|z_2\rangle\langle z_2|+\frac25|0,3\rangle\langle0,3|,
 \qquad
 z_2=\frac{3|00\rangle+\sqrt{21}|11\rangle
                       +\sqrt{21}|22\rangle+3|33\rangle}{10}.
\]
Here the local label $0,1,2,3$ encodes $00,01,10,11$ on two copies.
The marginal diagonals are exactly
$(49,21,21,9)/100$ and $(9,21,21,49)/100$.
Symmetry forces $t_*=1/2$, giving the exact regularized value
\[
 \boxed{\lim_{n\to\infty}C_n^{1/n}=\frac{\sqrt{21}}5
       \simeq0.91651514.}
\]
Consequently the ratio to independent one-copy optimization grows
exponentially at rate
\[
 \lim_{n\to\infty}\left(\frac{C_n}{C_1^n}\right)^{1/n}
       =\frac{\sqrt{21}}3>1.
\]
This is a gain in a prescribed-target coupling optimization. It is not a
claim that the fully entangled fraction of a fixed tensor-product state
equals $C_n$, or that $C_n$ approaches one: in this example $C_n$ tends
to zero despite the exponentially growing relative advantage.

Strict supermultiplicativity also does not imply monotonicity of the
per-copy quantities. Here
\[
 C_3=\frac{297+27\sqrt{21}}{1000},\qquad
 C_2^3-C_3^2=\frac{2349\sqrt{21}-4239}{500000}>0,
\]
so $C_3^{1/3}\simeq0.749321<C_2^{1/2}\simeq0.758258$.
The asymptotic limit remains exact; the parity effect below explains
these finite-copy oscillations.

\subsection{Exact parity-dependent prefactors for the symmetric binary family}
Let $a=(t,1-t)$, $b=(1-t,t)$, $p=(1/2,1/2)$ with $1/2<t<1$.
Write $r=\sqrt{(1-t)/t}\in(0,1)$ and $\lambda=2\sqrt{t(1-t)}$.
For $K_n\sim\operatorname{Binomial}(n,1/2)$ the exact formula becomes
\begin{equation}\label{eq:binary-expectation}
 C_n=\lambda^n\left[\mathbb E r^{|K_n-n/2|}\right]^2.
\end{equation}
The following sharper asymptotics distinguish even and odd copy numbers:
\begin{align}
 C_n&\sim\lambda^n\frac2{\pi n}
           \left(\frac{1+r}{1-r}\right)^2,
             &&n\to\infty\text{ through even integers},\label{eq:binary-even}\\
 C_n&\sim\lambda^n\frac2{\pi n}
           \left(\frac{2\sqrt r}{1-r}\right)^2,
             &&n\to\infty\text{ through odd integers}.\label{eq:binary-odd}
\end{align}
To prove them, divide each probability mass in the expectation by its
central binomial mass. For every fixed distance from the centre this
ratio tends to one, and it is always at most one. Dominated convergence
applies to the summable geometric weights. The limiting weight sums are
$1+2\sum_{j\ge1}r^j=(1+r)/(1-r)$ for even $n$, and
$2\sum_{j\ge0}r^{j+1/2}=2\sqrt r/(1-r)$ for odd $n$.
Stirling's formula gives central mass asymptotic
$\sqrt{2/(\pi n)}$, which proves both displayed equivalents.
The $n^{-1}$ prefactor shows why the exponential rate alone substantially
overestimates finite-copy capacities.

\paragraph{Attribution and verification.}
The coupling-capacity framework is due to Skalski, Todorov and Turowska
\cite{capacity2025}. The classical error-exponent method originates with
Chernoff and is reviewed, including geometric mixtures and unnormalized
divergence formulations, by Nielsen \cite{nielsen2022}. Neither that
method nor Stirling's formula is claimed as new. The additional result is
the exact collective quantum-coupling formula, its strict equality
criterion, explicit optimal states, and the resulting regularized value
and binary prefactors. A bounded search did not locate this specific
statement; that is not a proof of historical priority. The asymptotic
argument is an exact written proof, not a Lean-checked probability theorem.
